% AY2016 材料力学 〔3〕（転記）
% 出典: 04_材料力学/original/04_材料力学/2016/2016za.pdf
% 要約: 外径d, 長さLの円形断面はり（固定端A, 自由端B）。先端Bに長さL/2の剛体棒,
%       その先端に外力F。縦弾性係数E, 横弾性係数G。点Q: x=L-a, y=0, z=d/2。
%       (1) σx と 円周方向σy, (2) τxy, (3) Q点モール円・最大主応力,
%       (4) はりを〔2〕の薄肉円筒に置換時の σx,σy,τxy（Iz, Ip を用いる）。
% rem: 原本は3次元俯瞰図。軸構成（x:軸方向, y:上, z:横）を保った簡略立体図で再現する。

\section*{〔3〕}

外径 $d$, 長さ $L$ の円形断面はりの先端 B に長さ $L/2$ の剛体棒がある.
この剛体棒の先端に外力 $F$ が負荷されている. 縦弾性係数, 横弾性係数は $E$, $G$ とする.
（$Q$: $x=L-a$, $y=0$, $z=d/2$）

\begin{center}
\begin{tikzpicture}[>=Latex]
  % 固定壁（左, 斜めスラブ）
  \fill[pattern=north east lines] (0,-0.9) -- (0,1.3) -- (0.9,1.7) -- (0.9,-0.5) -- cycle;
  \draw (0,-0.9) -- (0,1.3) -- (0.9,1.7) -- (0.9,-0.5) -- cycle;
  % はり（円柱）本体: x方向（A:左 固定端 → B:右 自由端）
  \draw (0.6,0.4) -- (6.5,0.4);
  \draw (0.6,-0.4) -- (6.5,-0.4);
  \draw (6.5,0) ellipse (0.22 and 0.4);   % 自由端Bの断面
  % 軸（A点）
  \draw[->] (0.6,0) -- (1.9,0) node[above]{$x$};
  \draw[->] (0.6,0) -- (0.6,1.5) node[left]{$y$};
  \draw[->] (0.6,0) -- (1.3,0.65) node[right]{$z$};
  \node[below left] at (0.6,-0.2) {A};
  \node[below right] at (6.5,-0.3) {B};
  % 点 Q（上面, x=L-a）
  \fill (4.6,0.4) circle (1.8pt);
  \node[above] at (4.6,0.45) {Q};
  % 剛体棒（Bから上右へ, 長さL/2）と荷重F
  \draw[line width=1.4pt] (6.5,0.1) -- (7.6,1.0);
  \node at (7.3,0.45) {$L/2$};
  \draw[->,very thick] (7.6,2.1) -- (7.6,1.05) node[above]{$F$};
  % 寸法 L（はり長さ）
  \draw[<->] (0.6,-0.95) -- (6.5,-0.95) node[midway,fill=white]{$L$};
  % 寸法 a（Q→B）
  \draw[<->] (4.6,-0.55) -- (6.5,-0.55) node[midway,fill=white,font=\scriptsize]{$a$};
\end{tikzpicture}
\end{center}

\begin{enumerate}[label=(\arabic*)]
  \item $x$ 方向の応力 $\sigma_x$ と円周方向の $\sigma_y$ を求めよ.

  \item せん断応力 $\tau_{xy}$ を求めよ.

  \item Q 点におけるモール応力円を描け. そして, 最大主応力を求めよ.

  \item 円形断面はりを〔2〕のように薄肉円筒に置きかえる. このとき, $\sigma_x$, $\sigma_y$,
        $\tau_{xy}$ を求めよ. 断面二次モーメント, 断面二次極モーメントは $I_z$, $I_p$ とする.
\end{enumerate}
